\section{研究问题的细化及评价原则}
\subsection{为何把配置点选择单独作为研究对象}
在局部随机特征方法中，隐藏层参数、子域划分、配置点、损失权重和最小二乘求解器都能影响精度。如果同时调整所有因素，即使最终误差降低，也难以判断收益来自更好的表达空间、更多约束，还是更合适的选点。本文因此固定两个子域及隐藏特征，再考察点的数量和位置，把问题缩小为本科阶段可以独立验证的有限维研究。

这个缩小并没有消除所有交互。局部特征在不同位置可能具有不同的导数尺度，非均匀点又会改变观测矩阵，截断SVD随之改变有效空间。因此所谓固定网络并不意味着数值上保留的方向完全固定。本文报告奇异值、数值秩和物理误差，目的就是避免把所有变化都归因于残差指标。

界面问题还多了一种资源竞争：内部点主要约束各侧PDE，界面点同时约束值和通量。给界面更多物理位置会产生更多标量方程，而给内部更多点可能改善区域覆盖。对一种资源有利的安排，可能减少另一种资源。因此本课题的核心不是简单地证明残差越大越应多放点，而是检验这种直觉在预算、尺度与固定空间共同限制下是否成立。

\subsection{三个可检验的问题}
第一个问题是正确性：从模型、导数到矩阵的全部环节是否一致？这一问题由制造解一致性、非零跳跃、差分收敛和单元测试回答。没有通过这一层的算法比较没有解释基础；特别是高对比度时，小解误差也不能替代通量行的单独检验。

第二个问题是效果：在相同最终资源下，联合指标是否比静态点和仅PDE指标更有利？本文同时采用点预算与方程预算，并在每个设置内共享隐藏特征和随机种子。比较指标不仅包括相对解误差，还包括最大误差、独立PDE残差、两类界面误差及时间。不同指标不强制给出相同排序。

第三个问题是解释：如果联合方案无效或只有局部收益，哪些观测可以排除过强的解释？同配额固定随机对照可检验残差选点是否为消融收益的必要条件，宽度与范围扫描揭示表达空间和条件性变化，局部峰避免只在平滑低频模型上评价。它们仍不能构成完整因果证明，但能限制结论的范围。

\subsection{研究贡献的三个层次}
第一层是实现贡献，即把已有数学方法独立组织成可运行、可追溯的程序，并用明确模型验证。这一层要求代码与文稿一致，不能因为技术已有就省略复现细节，也不能把实现已有技术说成发明该技术。

第二层是实验与分析贡献，即在明示预算和参数下获得可重复的比较，解释观察到的取舍，保留与预期不一致的结果。对于本科论文，这种工作可以体现综合运用分析、代数、数值计算与概率统计知识的能力。其价值在于回答具体问题，而不是实验次数越多就自动意味着结论越强。

第三层才是方法改进贡献。只有相对于合理基线、在独立问题上重复验证效果，并排除额外资源或参数调优等解释后，才适合提出有限范围的改进主张。本文原计划探索这一层，但当前结果尚未达到相应证据门槛，因此最终定位以第二层为主。这是根据真实结果调整论断，不是放弃原来的研究问题。

\section{制造解研究的适用范围}
制造解允许精确计算真解、源项和跳跃，从而区分代码错误与近似误差。本文不依赖真实材料观测，也不需要训练数据标签；PDE及界面数据本身构成求解约束，解析真解只用于构造测试问题和独立误差评估，不进入残差选点。这样的设置适合验证数值算法，但与未知真解的真实应用仍有距离。

制造解的选择也会形成偏差。如果只使用非常平滑且与随机特征容易匹配的函数，配置点差异可能很小；如果使用超过当前固定空间分辨能力的尖峰，所有选点策略都可能精度不足。本文因此同时保留平滑、局部峰和复杂界面三类问题，而不把其中一类的排序外推到全部椭圆型问题。

系数对比度的方向需要区分。把$\beta=(1000,1)$换成$(1,1000)$不仅改变某个比值的写法，还改变高系数所在侧及制造解中系数依赖项。当前模型虽有规则几何，真解并非在所有参数下严格左右对称，因此正反向都实际运行。这比只报告单一方向更有利于检查尺度和组分配的影响。

需要注意，具有解析解不等于连续误差可以被机器精确积分。本文使用独立规则点或Gauss型积分近似范数，另做密度敏感性检查。图表中的数值是明确验证规则下的结果，不是经过区间算术认证的全域上下界。对于极小误差及高度振荡界面，积分规则和浮点求值本身也可能成为结果解释的一部分。

本论文最后形成的是可供逐项复核的计算研究。数学证明只覆盖已写明假设，数值结论只覆盖真实运行，尚未取得的原文和本届院系细则列为未决事项。这样的边界使读者能够知道哪些内容可以直接复算，哪些仍需要更强理论、作者代码或新的独立测试。
